% Einheit 2 von 4 – Terme zusammenfassen (bank/terme/e2.jsonl, zusammenbau v0.1)
%% TODO Zweigzeile Teil 1: was hier gelernt wird (Satz in Schülersprache aus dem Einheitstitel) – Einheit 2
\einheitenkopf[e2]{Einheit 2 von 4 · Terme zusammenfassen}
\zweigzeile{ab Kl. 7, je nach Buch bis Kl. 8 · P10}
\setcounter{aufgabe}{10}

%% TODO Titel: Ich-kann-Satz fehlt in der Bank (Platzhalter: Kettenname) – Nr. 11
%% TODO Anweisung: ein Satz über der ersten Teilaufgabe fehlt in der Bank – Nr. 11
\begin{aufgabe}{Gleichartige Glieder erkennen}
\begin{teile}
\teil Unterstreiche, was zusammengehört: $4y$, $7$, $9y$
\end{teile}
\end{aufgabe}

%% TODO Titel: Ich-kann-Satz fehlt in der Bank (Platzhalter: Kettenname) – Nr. 12
%% TODO Anweisung: ein Satz über der ersten Teilaufgabe fehlt in der Bank – Nr. 12
\begin{aufgabe}{Vorzahl lesen}
\begin{teile}
\teil $a$ – Vorzahl?
\end{teile}
\end{aufgabe}

%% TODO Titel: Ich-kann-Satz fehlt in der Bank (Platzhalter: Kettenname) – Nr. 13
%% TODO Anweisung: ein Satz über der ersten Teilaufgabe fehlt in der Bank – Nr. 13
\begin{aufgabe}{Vorzeichen als Teil des Gliedes}
\begin{teile}
\teil Glieder mit Vorzeichen: $5a - 2b + 7$
\end{teile}
\end{aufgabe}

%% TODO \verfahren{…}: Verfahrensüberschrift in Sachsprache fehlt in der Bank (2.3 b)
%% TODO Titel: Ich-kann-Satz fehlt in der Bank (Platzhalter: Kettenname) – Nr. 14
%% TODO Anweisung: ein Satz über der ersten Teilaufgabe fehlt in der Bank – Nr. 14
\begin{aufgabe}{Zusammenfassen}
%% TODO Darstellung für schwach fehlt in der Bank (terme-e2-k4-s1-v1; Abschnitt „Für schwache Schüler“)
\swz{Fasse zusammen: $2x + 5x$}{}
%% TODO Darstellung für schwach fehlt in der Bank (terme-e2-k4-s1-v2; Abschnitt „Für schwache Schüler“)
\swz{Fasse zusammen: $6a + 3a$}{}
%% TODO Darstellung für schwach fehlt in der Bank (terme-e2-k4-s1-v3; Abschnitt „Für schwache Schüler“)
\swz{Fasse zusammen: $4y + 4y$}{}
%% TODO Darstellung für schwach fehlt in der Bank (terme-e2-k4-s1-v4; Abschnitt „Für schwache Schüler“)
\swz{Fasse zusammen: $7b + 2b$}{}
%% TODO Darstellung für schwach fehlt in der Bank (terme-e2-k4-s1-v5; Abschnitt „Für schwache Schüler“)
\swz{Fasse zusammen: $5z + 6z$}{}
\swfrage{Was bleibt gleich, was ändert sich?}
%% TODO Darstellung für schwach fehlt in der Bank (terme-e2-k4-s2-v1; Abschnitt „Für schwache Schüler“)
\swz{Fasse zusammen: $9x - 4x$}{}
%% TODO Darstellung für schwach fehlt in der Bank (terme-e2-k4-s3-v1; Abschnitt „Für schwache Schüler“)
\swz{Fasse zusammen: $2x + 4x + 3x$}{}
%% TODO Darstellung für schwach fehlt in der Bank (terme-e2-k4-s4-v1; Abschnitt „Für schwache Schüler“)
\swz{Fasse zusammen: $5x + 3 + 2x$}{}
%% TODO Darstellung für schwach fehlt in der Bank (terme-e2-k4-s5-v1; Abschnitt „Für schwache Schüler“)
\swz{Fasse zusammen: $4a + 3b + 2a$}{}
%% TODO Darstellung für schwach fehlt in der Bank (terme-e2-k4-s6-v1; Abschnitt „Für schwache Schüler“)
\swz{Fasse zusammen: $y + 4y$}{}
%% TODO Darstellung für schwach fehlt in der Bank (terme-e2-k4-s7-v1; Abschnitt „Für schwache Schüler“)
\swz{Fasse zusammen: $2x - 7x$}{}
%% TODO Darstellung für schwach fehlt in der Bank (terme-e2-k4-s8-v1; Abschnitt „Für schwache Schüler“)
\swz{Fasse zusammen: $-3x + 8x$}{}
%% TODO Darstellung für schwach fehlt in der Bank (terme-e2-k4-s9-v1; Abschnitt „Für schwache Schüler“)
\swz{Fasse zusammen: $4x^2 + 5x + 2x^2$}{}
%% TODO Darstellung für schwach fehlt in der Bank (terme-e2-k4-s10-v1; Abschnitt „Für schwache Schüler“)
\swz{Fasse zusammen: $1{,}5x + 3{,}5x$}{}
\swz[3]{Vereinfache den Term $8a - 2a^2 - 5a$ und berechne seinen Wert für $a = 3$. \hfill (P10 2025 GYM)}{}
\end{aufgabe}

%% TODO \verfahren{…}: Verfahrensüberschrift in Sachsprache fehlt in der Bank (2.3 b)
%% TODO Titel: Ich-kann-Satz fehlt in der Bank (Platzhalter: Kettenname) – Nr. 15
%% TODO Anweisung: ein Satz über der ersten Teilaufgabe fehlt in der Bank – Nr. 15
\begin{aufgabe}{Malnehmen}
%% TODO Darstellung für schwach fehlt in der Bank (terme-e2-k5-s1-v1; Abschnitt „Für schwache Schüler“)
\swz{Multipliziere: $5 \cdot 2x$}{}
%% TODO Darstellung für schwach fehlt in der Bank (terme-e2-k5-s1-v2; Abschnitt „Für schwache Schüler“)
\swz{Multipliziere: $4 \cdot 6a$}{}
%% TODO Darstellung für schwach fehlt in der Bank (terme-e2-k5-s1-v3; Abschnitt „Für schwache Schüler“)
\swz{Multipliziere: $7 \cdot 3y$}{}
%% TODO Darstellung für schwach fehlt in der Bank (terme-e2-k5-s1-v4; Abschnitt „Für schwache Schüler“)
\swz{Multipliziere: $2 \cdot 9b$}{}
%% TODO Darstellung für schwach fehlt in der Bank (terme-e2-k5-s1-v5; Abschnitt „Für schwache Schüler“)
\swz{Multipliziere: $6 \cdot 5z$}{}
\swfrage{Was bleibt gleich, was ändert sich?}
%% TODO Darstellung für schwach fehlt in der Bank (terme-e2-k5-s2-v1; Abschnitt „Für schwache Schüler“)
\swz{Multipliziere: $4x \cdot 5$}{}
%% TODO Darstellung für schwach fehlt in der Bank (terme-e2-k5-s3-v1; Abschnitt „Für schwache Schüler“)
\swz{Multipliziere: $4y \cdot 2y$}{}
%% TODO Darstellung für schwach fehlt in der Bank (terme-e2-k5-s4-v1; Abschnitt „Für schwache Schüler“)
\swz{Multipliziere: $(-4) \cdot 5x$}{}
%% TODO Darstellung für schwach fehlt in der Bank (terme-e2-k5-s5-v1; Abschnitt „Für schwache Schüler“)
\swz{Multipliziere: $4a \cdot 5b$}{}
%% TODO Darstellung für schwach fehlt in der Bank (terme-e2-k5-s6-v1; Abschnitt „Für schwache Schüler“)
\swz{Multipliziere: $2 \cdot 5x \cdot 3$}{}
\swz[3]{Multipliziere: $(-6a) \cdot 2b$}{}
\end{aufgabe}

%% TODO Titel: Ich-kann-Satz fehlt in der Bank (Platzhalter: Kettenname) – Nr. 16
%% TODO Anweisung: ein Satz über der ersten Teilaufgabe fehlt in der Bank – Nr. 16
\begin{aufgabe}{Zusammenfassen – Fehler finden, Begründen, Darstellungswechsel, Anwendung}
\swz[3]{Finde den Fehler und rechne richtig: \rechnung{5a + 3 + 2a &= 10a}}{}
\swfrage{Kann man $4x + 4y$ zusammenfassen? Begründe.}
\swz{Umfang des Rechtecks (Maße in cm) – Term, zusammengefasst?}{\viereck[seiten={3x,2,3x,2}]{(0,0)}{(6,0)}{(6,2)}{(0,2)}}
\swz[3]{Ein Kiosk verkauft Brötchen zu je $p$ Euro: am Montag 45, am Dienstag 38, am Mittwoch 52 Stück. Stelle einen Term für die Einnahmen der drei Tage auf und fasse zusammen. Wie viel nimmt der Kiosk bei $p = 0{,}40$ ein?}{}
\end{aufgabe}

\uebersichtskasten{Zusammenfassen: Nur gleiche Variablen mit gleicher Hochzahl darf man zusammenfassen – Vorzahlen addieren, Variable bleibt. \\ 3x + 5x = 8x \quad 7a − 2a + 4 = 5a + 4 \quad x² + 3x bleibt so \\ Malnehmen: Zahlen mal Zahlen, Variablen mal Variablen. \\ 3 · 4x = 12x \quad 2x · 3x = 6x² \quad 3a · 2b = 6ab}
